\chapter{Conclusion and Future Work}

\section{Conclusion}

This thesis designed and implemented a hybrid Earth Observation metadata infrastructure for
research and education in the Kvarken Space Center context. The central design decision was to
treat reliability and reproducibility as properties of the complete path from provider response to
regional query, rather than as features added after ingestion. The resulting system combines
typed boundary validation, injectable provider transports, bounded asynchronous scheduling,
content-addressed provenance, a WAL-enabled SQLite spatial catalog, metadata quality profiling,
an educational HTTP interface, and a small raster/NDVI analytical slice.

\section{Answers to the research questions}

\subsection{Research question 1: metadata normalization and provenance}

Heterogeneous STAC and CDSE metadata can be normalized by validating provider JSON at the
transport boundary and transforming only accepted Items into a provider-neutral frozen
\texttt{EOScene} dataclass. The model retains identity, platform, acquisition time, geometry,
coordinate reference, quality metadata, and asset links without exposing provider-specific
extensions to downstream code. Provenance is preserved independently: canonical raw payload bytes
receive a SHA-256 content address, while an append-only JSON Lines manifest records source,
timestamp, digest, and retry context. This pairing makes the normalized catalog convenient to
query without discarding the evidence needed to reconstruct an ingestion decision.

\subsection{Research question 2: bounded recovery}

Asynchronous ingestion remains bounded by placing an \texttt{asyncio.Semaphore} around the
expensive transformation and persistence path. If the configured worker limit is \(k\), no more
than \(k\) tasks are active at once. Recovery is explicit and selective: timeouts, HTTP 429
responses, and token expiry use bounded exponential backoff and refresh; malformed payloads and
invalid token structures are surfaced as permanent errors. The 53-test offline suite demonstrates
these behaviors with injectable transports rather than relying on nondeterministic provider
incidents.

\subsection{Research question 3: regional catalog behavior}

A dependency-free SQLite catalog can support the required regional, temporal, quality, and
platform queries when rectangular predicates are indexed before deterministic polygon checks.
The Kvarken EPSG:4326 bounding box provides a clear contract, while boundary-touching footprints
are retained by inclusive intersection semantics. WAL mode permits readers to observe committed
snapshots while writers append transactions, and contention tests verify valid intermediate
states and durability after reopening. This is appropriate for an embedded educational and
research catalog, although it is not a substitute for a distributed database under multi-node
load.

\subsection{Research question 4: reproducible empirical evidence}

An offline baseline is reproducible when its fixtures, concurrency setting, report schema, source
code, and validation commands are versioned together. The frozen 256-scene artifact records
220.819653 items per second, zero failures and retries, complete catalog and asset coverage,
0.106130 MB of metadata payload, and 9.117 ms spatial query latency in the measured environment.
The report is intentionally additive: authenticated live runs produce dated artifacts without
overwriting this comparison point. Thus the baseline supports regression analysis while clearly
separating synthetic evidence from future provider observations.

\section{Scientific and architectural contributions}

The main contribution is a coherent set of small design mechanisms that can be inspected and
tested independently:

\begin{itemize}
\item A dependency-free runtime demonstrates that asynchronous EO metadata ingestion, OAuth2
boundary handling, SQLite spatial indexing, HTTP querying, provenance, reporting, and NDVI
calculation do not require a heavyweight framework for a research baseline.
\item Standard-library asynchronous scheduling makes the concurrency contract explicit and
testable, preventing an apparently parallel pipeline from creating unbounded provider or local
work.
\item SHA-256 content addressing and append-only manifests separate immutable source evidence from
mutable query indexes and make duplicate or missing provenance detectable.
\item SQLite WAL persistence provides a compact indexed regional catalog with reader/writer
isolation, transactional upserts, and durability checks suitable for a single-node research
deployment.
\item The frozen benchmark, deterministic reports, fault-injection tests, and health command turn
implementation claims into a repeatable empirical protocol rather than an unversioned demo.
\end{itemize}

The contribution is consequently methodological as well as software-oriented: it defines a
traceable path for moving from an external EO standard to an operationally meaningful, teachable
research artifact.

\section{Future work}

\subsection{Authenticated continuous CDSE ingestion}

The current CDSE adapter supports client-credentials authorization and bounded recovery, but CI
uses offline fixtures. A next phase should run a scheduled continuous ingestion service with
rotating credentials, provider-aware rate budgets, checkpointed pagination, and monitoring for
schema drift. Each execution should retain the existing dated report and provenance contracts.
Comparisons should quantify provider latency, response distributions, token refresh frequency, and
retry cost against the frozen synthetic baseline.

\subsection{Raster and GeoTIFF array operations}

The raster slice currently validates HTTP range transport and computes NDVI over aligned numerical
values. Future work should add a carefully bounded GeoTIFF reader or an optional raster
dependency, windowed reads, nodata masks, scale/offset handling, and alignment checks across red
and near-infrared bands. Evaluation should measure memory use and numerical agreement on
representative Sentinel-2 windows without weakening the metadata-only offline path.

\subsection{Distributed spatial catalogs}

SQLite WAL is an effective single-node boundary, but multi-node ingestion would require a
different ownership and consistency model. A distributed extension could partition scenes by time
or region, use a service-backed spatial index, and preserve content-addressed provenance in shared
object storage. The existing typed catalog protocol and deterministic regional predicates should
remain the compatibility contract while benchmarks compare write contention, query tail latency,
recovery, and operational cost.

\subsection{Operational deployment at Kvarken Space Center}

Deployment research should connect the pipeline to an operational scheduling environment at
Kvarken Space Center. This includes secret management, durable raw/staged storage, observability,
backup and restore drills, access control, data-retention policy, and an educational dashboard
over the read-only query API. A staged rollout should retain offline replay tests and the frozen
256-scene artifact as release gates so that operational convenience does not erase research
traceability.

\section{Closing statement}

The released baseline establishes a practical starting point rather than a finished planetary
data platform. Its value is the explicitness of its contracts: what is accepted, what is retried,
what is stored, how a region is defined, and how a result is measured. Future provider,
processing, and deployment work can therefore extend the system while preserving the evidence and
comparability required for a credible Master's thesis.
